\subsection{由等分容器过渡到任意体积分区}
“左、右各一半”只是几何概率的一种特殊情形。粒子处在某区域的概率，先由允许位置的测度决定；再用独立试验规则计算粒子数的分布。固定总数并不表示各区域粒子数固定。

\begin{origexamplebox}{体积比为四比一的两区域粒子数}\label{ch18:ex:teaching-asymmetric-binomial}
均匀平衡的稀薄理想气体含 $N$ 个粒子。把容器在几何上分成体积为 $4V/5$ 和 $V/5$ 的两区域，并无阻止粒子穿过的实体隔板。忽略粒子位置间的关联，求较小区域粒子数 $n$ 的分布、平均值和标准差，并解释宏观密度为何仍近似稳定。
\end{origexamplebox}
\begin{xesolution}
每个粒子在小区域的概率为 $q=1/5$。选定哪 $n$ 个粒子进入小区域有 $\binom Nn$ 种组合，每种组合的概率为 $q^n(1-q)^{N-n}$，所以
\[
 P(n)=\binom Nn(1/5)^n(4/5)^{N-n},\qquad 0\le n\le N.
\]
二项式定理保证全部概率之和为1。令 $X_i$ 表示第 $i$ 个粒子是否在小区域，$X_i=1$ 的概率为 $q$。由于 $X_i^2=X_i$，有 $\langle X_i\rangle=q$、$\operatorname{Var}(X_i)=q(1-q)$。独立性消去协方差，故
\[
 \langle n\rangle=N/5,\qquad \operatorname{Var}(n)=4N/25,
 \qquad \sigma_n=2\sqrt N/5,\qquad
 \frac{\sigma_n}{\langle n\rangle}=\frac2{\sqrt N}.
\]
粒子数的绝对起伏随 $N$ 增大，却比平均数增长得慢，所以相对密度起伏变小。大区域粒子数为 $N-n$，与小区域完全反相关；不能把两区域总粒子数的方差写成两个方差直接相加，否则会误判固定的 $N$ 也在涨落。这里的临时粒子编号只用于稀薄经典模型的位置计数，不是为全同粒子添加新的可观测身份。
\end{xesolution}
